\section{Tensors and Tensor Fields}

\subsection{Tensor Algebra}

We generalize bilinear forms to multilinear forms.

\begin{definition}[Multilinear map and multilinear function]
Let $V_1, \dots, V_k$ and $W$ be real vector spaces. A map
\[
f: V_1 \times \dots \times V_k \to W
\]
is \emph{multilinear} if it is linear in each component. For example,
\[
f(a v_1 + b w_1, v_2, \dots, v_k)
= a f(v_1, v_2, \dots, v_k) + b f(w_1, v_2, \dots, v_k)
\]
for all $a,b\in\mathbb{R}$ and $v_1,w_1\in V_1$. If $W=\mathbb{R}$, $f$ is called a \emph{multilinear function}.
\label{def:multilinear}
\end{definition}

\begin{definition}[$(s,t)$-type tensor]
Let $V$ be an $n$-dimensional real vector space and $V^*$ its dual.

For $s,t \in \mathbb{Z}_{\ge 0}$, a multilinear map
\[
T : (V^*)^s \times V^t \to \mathbb{R},
\]
is called a \emph{tensor of type} $(s,t)$. The set of all such tensors is denoted by $T_t^s(V)$; it is a real vector space under pointwise addition and scalar multiplication.
\end{definition}

\begin{remark}
Since the evaluation map $V \to V^{**}$
$v^{**}(f):= f(v)$ is a linear isomorphism,
it is sufficient to consider only $V^*$ and $V$ in definition \ref{def:multilinear} 
\end{remark}

\begin{problem}
For a multilinear map $T: (V^*)^s \times V^t \to \mathbb{R}$, define $\hat{T} \in T_t^{s+1}(V)$ by
\[
\hat{T}(\omega^1, \dots, \omega^{s+1}, v_1, \dots, v_t)
:= \omega^{s+1}\bigl(T(\omega^1, \dots, \omega^s, v_1, \dots, v_t)\bigr),
\]
where $\omega^1, \dots, \omega^{s+1} \in V^*$ and $v_1, \dots, v_t \in V$.
Show that the map $T \mapsto \hat{T}$ is a linear isomorphism.
\end{problem}

\begin{definition}[Tensor product]
For $T \in T_t^s(V)$ and $S \in T_{t'}^{s'}(V)$, define $T \otimes S \in T_{t+t'}^{s+s'}(V)$ by
\[
\begin{aligned}
(T \otimes S)&(\omega^1, \dots, \omega^{s+s'}, v_1, \dots, v_{t+t'})\\
&:=\; T(\omega^1, \dots, \omega^s, v_1, \dots, v_t) 
\times\; S(\omega^{s+1}, \dots, \omega^{s+s'}, v_{t+1}, \dots, v_{t+t'}),
\end{aligned}
\]
where $\omega^1, \dots, \omega^{s+s'} \in V^*$ and $v_1, \dots, v_{t+t'} \in V$.
\end{definition}

\begin{remark}
For any tensors $T, S, R$ and $a \in \mathbb{R}$, we have:
\begin{enumerate}
\item $\left(T \otimes S \right) \otimes R = T \otimes \left(S \otimes R\right)$.
\item $\left(aT\right) \otimes S = T \otimes \left(a S \right) = a \left(T \otimes S\right)$.
\end{enumerate}
These are denoted as $T \otimes S \otimes R$ and $a T \otimes S$, respectively.
Furthermore, $S$ and $R$ have the same type, the distributive laws hold :
\begin{enumerate}
    \item $T \otimes (S + R) = T \otimes S + T \otimes R$.
    \item $(S+R) \otimes T = S \otimes T + R \otimes T$.
\end{enumerate}
\end{remark}

\begin{remark}
In general, $S \otimes T \neq T \otimes S$.
\end{remark}

Let $\{e_i \}$ be a basis for $V$ and $\{e^i\}$ be its dual.
For $T \in T_t^s(V)$, the $n^{s+t}$ real numbers
\[
T^{i_1 \dots i_s}_{j_1 \dots j_t}
:= T(e^{i_1}, \dots, e^{i_s}, e_{j_1}, \dots, e_{j_t}),
\]
are called the \emph{coefficients} (or \emph{components}) of $T$.
For $T \in T_t^s(V)$ and $S \in T_{t'}^{s'}(V)$, the coefficients of $T \otimes S$ are given by
\[
 (T\otimes S)^{i_1 \dots i_{s+s'}}_{j_1 \dots j_{t+t'} }
 = T^{i_1 \dots i_s}_{j_1 \dots j_t} \; S^{i_{s+1} \dots i_{s+s'}}_{j_{t+1} \dots j_{t+t'}}.
\]

\begin{problem}
Show that the set of $n^{s+t}$ tensors
\[
e_{i_1} \otimes \dots \otimes e_{i_s} \otimes e^{j_1} \otimes \dots \otimes e^{j_t},
\]
where $i_1, \dots, i_s, j_1, \dots, j_t = 1, \dots, n$ 
forms a basis for $T_t^s(V)$; and that for any $T \in T_t^s(V)$,
\[
T = \sum_{i_1, \dots, i_s, j_1, \dots, j_t} T_{j_1 \dots j_t}^{i_1 \dots i_s} e_{i_1} \otimes \dots \otimes e_{i_s} \otimes e^{j_1} \otimes \dots \otimes e^{j_t}.
\]
In particular, show that $\dim T_t^s(V) = n^{s+t}$. We often write 
\[
T_t^s(V) = V \otimes \dots \otimes V \otimes V^* \otimes \dots \otimes V^*.
\]
\end{problem}

\subsection{Exterior Algebra}

\begin{definition}[p-form]
    Let V be a real vector space and $p \in \mathbb{N}$.
    An element of $T_p^0(V)$ is called a \emph{p-form}.
\end{definition}
\begin{remark}
A p-form $\omega$ is \emph{skew symmetric}. That is,
\[
\omega(v_{\sigma(1)}, \dots, v_{\sigma(p)}) = \text{sgn}(\sigma) \omega(v_1, \dots, v_p),
\]
for any $v_1, \dots, v_p \in V$ and $\sigma \in S_p$.
Furthermore, 
$\Lambda^p V^* := \{ \text{skew symmetric p-forms} \}$ is a vector space.
\end{remark}

\begin{definition}[Exterior product]
    For $\alpha \in \Lambda^p(V^*)$ and $\beta \in \Lambda^q(V^*)$, 
    the \emph{exterior product} $\alpha \wedge \beta$ is defined by 
    \[
    \alpha \wedge \beta := \frac{1}{p! q!} \sum_{\sigma \in S_{p+q}} \text{sgn}(\sigma) \alpha \otimes \beta \left(v_{\sigma(1)}, \dots, v_{\sigma(p+q)}\right).
    \]
\end{definition}

\begin{proposition}\label{prop:exterior_product}
    Let $\alpha, \alpha^1, \alpha^2 \in \Lambda^p(V^*)$ and $\beta, \beta^1, \beta^2 \in \Lambda^q(V^*)$.
    \begin{enumerate}
        \item $\left( a \alpha^1 + b \alpha^2 \right) \wedge \beta = a \alpha^1 \wedge \beta + b \alpha^2 \wedge \beta$.
        \item $\alpha \wedge \left( a \beta^1 + b \beta^2 \right) = a \alpha \wedge \beta^1 + b \alpha \wedge \beta^2$.
        \item $\left(\alpha \wedge \beta \right) \wedge \gamma = \alpha \wedge \left( \beta \wedge \gamma \right)$.
        \item $\beta \wedge \alpha = (-1)^{pq} \alpha \wedge \beta$.
    \end{enumerate}
\end{proposition}

\subsection{Tensor Fields and differential forms on manifolds}

Let M be an n-dimensional manifold an $\left( V, \phi \right)$ a chart with
$\phi = (x^1, \dots, x^n)$. At each $p \in U$, applying \autoref{prop:exterior_product} with $V = T_pM$, yields
\[
\frac{\partial}{\partial x^{i1}} \otimes \dots \otimes \frac{\partial}{\partial x^{is}} \otimes dx^{j1} \otimes \dots \otimes dx^{jt},
\]
which forms a basis for $T_t^s(T_pM)$. Any $T_p \in T_t^s (T_pM)$ can be written as,
\[
T_p = \sum_{i_1, \dots, i_s, j_1, \dots, j_t} T_{j_1 \dots j_t}^{i_1 \dots i_s} \frac{\partial}{\partial x^{i_1}} \otimes \dots \otimes \frac{\partial}{\partial x^{i_s}} \otimes dx^{j_1} \otimes \dots \otimes dx^{j_t}.
\]

\begin{definition}[$(s,t)$-type tensor field]
A map $T: M \owns p \to T_p \in T_t^s(T_pM)$ is a\emph{(s,t)-type tensor field}. $T$ is $C^{\infty}$ if all its coefficients 
$T_{j_1 \dots j_t}^{i_1 \dots i_s}$ are $C^{\infty}$ functions on any chart $(V, \phi)$.
\end{definition}

\begin{remark}
It suffices to check the $C^{\infty}$ property on a single atlas.
\end{remark}




